 % ==================== 1.2. ENFERMEDADES NEURODEGENERATICAS. VISIÓN GENERAL. ====================

 % #TODO AÑADIR CITAS EN ESTA SECCIÓN

\section{Enfermedades neurodegenerativas. Visión general}\label{sec:vision_general}

Las enfermedades neurodegenerativas engloban un conjunto de trastornos definidos por la degradación progresiva del sistema nervioso central, induciendo un deterioro gradual e irreversible de las capacidades motoras, cognitivas y lingüísticas del paciente. Aunque cada entidad nosológica exhibe un fenotipo clínico particular, se reconocen mecanismos fisiopatológicos compartidos, tales como la agregación proteica anómala, el estrés oxidativo, la neuroinflamación sostenida y la apoptosis neuronal selectiva en regiones encefálicas específicas \cite{abrilcarreres2024}. En la actualidad, este conjunto de patologías representa una causa primaria de discapacidad en la población geriátrica, previéndose un incremento exponencial de su prevalencia como consecuencia directa del envejecimiento demográfico global \cite{avance2023}. 

Dentro de este espectro clínico, la Enfermedad de Parkinson, la Enfermedad de Alzheimer y la Esclerosis Lateral Amiotrófica presentan un interés investigativo prioritario para el análisis acústico biomédico.

La \textbf{Enfermedad de Parkinson (EP)} se define como una patología crónica originada por la degeneración de las neuronas dopaminérgicas ubicadas en la sustancia negra mesencefálica. Clásicamente diagnosticada a través de manifestaciones motoras (temblor en reposo, bradicinesia, rigidez articular), los pacientes desarrollan con alta frecuencia una sintomatología fonatoria específica conocida como disartria hipocinética \cite{romero2018}. Se registra habitualmente hipofonía, pérdida de modulación frecuencial, imprecisión articulatoria y un déficit notable en el control del soporte respiratorio. Frecuentemente, estas anomalías se manifiestan en estadios subclínicos y preceden a la sintomatología motora apendicular convencional, posicionando al procesamiento digital del habla como una técnica notablemente sensible para su detección temprana \cite{mendes2024}.

Por otra parte, la \textbf{Enfermedad de Alzheimer (EA)} constituye el subtipo de demencia con mayor incidencia epidemiológica, caracterizándose por un declive progresivo de las funciones cognitivas superiores. Las alteraciones comunicativas asociadas a esta demencia exhiben una naturaleza fundamentalmente lingüística y prosódica \cite{garcia-gutierrezUnveilingSoundCognitive2024}. Se observa un incremento en la duración y recurrencia de las pausas silenciosas, una pérdida de afluencia discursiva, restricción del acervo léxico, simplificación sintáctica y un aplanamiento de la curva entonativa. Estas modificaciones deterioran tanto la estructura acústica de la emisión como la organización semántica del discurso espontáneo, permitiendo obtener inferencias clínicas precisas sobre el estado neurológico del individuo.

Respecto a la \textbf{Esclerosis Lateral Amiotrófica (ELA)}, esta enfermedad provoca la degeneración progresiva de las motoneuronas superiores e inferiores, derivando en una debilidad muscular incapacitante. Cuando la degeneración compromete inicialmente la región bulbar (estructura neuroanatómica responsable del control de los músculos implicados en la respiración, la fonación y la articulación) se manifiesta clínicamente mediante disartria espástica o flácida, inestabilidad en la intensidad fonatoria y fatiga respiratoria prematura \cite{romero2018}. La evaluación paramétrica de estas características acústicas facilita la identificación temprana del daño neuromuscular bulbar y habilita la monitorización no invasiva de la progresión de la enfermedad.

La práctica clínica habitual fundamenta el diagnóstico de estas patologías en la aplicación de baterías neuropsicológicas y la realización de pruebas de neuroimagen. Si bien estas metodologías convencionales poseen validez clínica, presentan restricciones significativas en términos de infraestructura, coste económico y sensibilidad en fases preclínicas. Ante este escenario, el modelado computacional de la voz mediante Inteligencia Artificial se propone como un sistema de apoyo a la decisión clínica complementario, ofreciendo una alternativa metodológica objetiva, accesible y con alta capacidad de escalabilidad poblacional \cite{moro-velazquezNewToolsDifferential2021}.